% !PRINTED 227
\weberpage{227}
%<<CONT
die Gruppennatur nachgewiesen und die Ikosa\"edergruppe gebildet ist.

Aus der ersten Formel (25) ergiebt sich noch
\[
\psi = \chi\Theta^{\mp1}\chi\Theta^{\pm1}\chi\Theta^{\mp1},
\]
woraus man schliesst, dass $\psi$ aus $\chi$ und $\Theta$ abgeleitet werden kann, dass also $\chi$ und $\Theta$ allein schon als \gesperrt{erzeugende Substitutionen der Ikosa\"edergruppe} betrachtet werden können.

Wollen wir die Substitution $\chi$ mit der Determinante 1 darstellen, so beachten wir die Relation
\[
-\omega^2-1 = \omega-2 = (\varepsilon^2-\varepsilon^{-2})^2,
\]
und erhalten, wenn wir noch $\varepsilon+\varepsilon^{-1}$ für $\omega$ einsetzen:
\weq{26}{\chi = \begin{pmatrix} \dfrac{1}{\varepsilon-\varepsilon^{-1}}, & \dfrac{1}{\varepsilon^2-\varepsilon^{-2}} \\[2mm] \dfrac{1}{\varepsilon^2-\varepsilon^{-2}}, & \dfrac{-1}{\varepsilon-\varepsilon^{-1}} \end{pmatrix}.}

\S.\,59.

\gesperrt{Die Theiler der Ikosa\"edergruppe.}

Die Theiler der Ikosa\"edergruppe müssen unter den niedrigeren Poly\"edergruppen gesucht werden. Darunter finden sich zunächst die cyklischen Gruppen $C_n$, deren jede aus der Periode einer der Ikosa\"edersubstitutionen besteht. Die Abzählung der Anzahl dieser Gruppen ergiebt sich sofort aus der Zusammenstellung der Substitutionen nach ihren Graden, wie wir sie im vorigen Paragraphen gegeben haben; nämlich:
\[
\begin{aligned}
&15 \text{ Gruppen } C_2, \\
&10 \quad \text{„} \quad C_3, \\
&\phantom{0}6 \quad \text{„} \quad C_5.
\end{aligned}
\]
Es sind ferner in der Ikosa\"edergruppe Di\"edergruppen $D_2$, $D_3$, $D_5$ enthalten, als deren Repräsentanten wir folgende aufstellen:
\[
\begin{aligned}
D_2 &: 1, \psi, \chi, \psi\chi, \\
D_3 &: 1, \chi\Theta^2, \Theta^{-2}\chi, \psi\chi, \psi\Theta^2, \Theta^{-2}\psi\chi\Theta^2, \\
D_5 &: \Theta^r, \psi\Theta^r, \quad r=0,1,2,3,4,5.
\end{aligned}
\]
% [sic: Weber prints r=0,1,2,3,4,5; since \Theta^5=1, the last value duplicates r=0.]
